\scanpage{6}\mspage{5}
\deleted{une figure a}

Pour \deleted{toute} figure $F$, $\operatorname{Multomb}(F)$ \inserted{est appelé} \deleted{\unsure{de|la}} la \uline{grande} \uline{multiombre} de $F$, ou la $\mathcal{M}$-multiombre de $\mathbb{F}$. \deleted{On} C'est en fait une figure ensembliste dans $\mathcal{M}$\inserted{indexée fidèlement par $\widetilde F$ (le déployement de $F$)}, et dont la connaissance implique celle de $F$.

\inserted{Jusqu'à présent il n'était pas question de $\mathcal{L}$.} Utilisant \deleted{d'autre part} \inserted{\unsure{maintent|maintnt|maintenant}} l'application
\[
\mathcal{L}\hookrightarrow\mathcal{M},
\]
et les applications images inverses correspondantes
\[
\begin{aligned}
\unsure{\mathcal{P}|\mathscr{P}|\mathfrak{P}}(\mathcal{M})&\longrightarrow\mathcal{P}(\mathcal{L})\\
\operatorname{Fig}(\mathcal{M})&\rightharpoonup\operatorname{Fig}(\mathcal{L})
\end{aligned}
\]
on trouve, en composant avec 1.6 a)b)
\[
\left\{\begin{aligned}
\text{a)}\quad \unsure{\mathcal{F}|\mathfrak{F}}&\mathrel{\deleted{\lhook}\!\longrightarrow}\mathcal{P}(\mathcal{L})&
 F&\longmapsto\operatorname{omb}(F)\\
\text{b)}\quad \mathcal{F}&\mathrel{\deleted{\lhook}\!\longrightarrow}\operatorname{Fig}(\mathcal{L})&
 F&\longmapsto\operatorname{mul}\unsure{t}\operatorname{omb}(F),
\end{aligned}\right.\tag*{(1.9)}
\]
% The t in this occurrence of multomb is heavily retraced or overwritten.
% Its live reading is t; no distinct previous character can be separated securely.
% The manuscript places the label (1.9) to the left of the brace.
dont la première n'a aucune raison d'être injective, la seconde par contre a une tendance à l'être $\times$ \deleted{on ne rien} \inserted{On} appelle $\operatorname{omb}(F)$ la \uline{petite ombre} de $F$, ou simplement son ombre quand une confusion semble exclue, et de $\overline m$ pour $\operatorname{multomb}$, la (\uline{petite}) \uline{multiombre de la} figure $F$.
